\documentclass[10pt,a4paper]{amsart}
\usepackage[margin=19mm]{geometry}
\usepackage{amsmath,amssymb,mathtools}
\usepackage[scr=rsfs,cal=euler]{mathalfa}
\usepackage{xfrac}
\usepackage[unicode,hidelinks,bookmarks=false]{hyperref}
\newcommand{\ie}{i.e.\ }
\newcommand{\cf}{cf.\ }
\newcommand{\resp}{respectively\ }
\newcommand{\Subsection}{\subsection}
\providecommand{\nobreakdash}{\nobreak\hbox{-}}
\setlength{\emergencystretch}{2em}
\multlinegap0pt
\usepackage{amssymb}
\usepackage[shortlabels]{enumitem}
\usepackage{tikz}
\usepackage{bbm}
\usepackage{multirow}
\usepackage{xcolor}
\usepackage{booktabs}
\usepackage{algorithm}\usepackage{algpseudocode}
\usepackage{dsfont}
\usepackage{float}
\usepackage{caption}
\newtheorem{lemma}{Lemma}
\newcommand{\D}{\mathrm{d}}
\newcommand{\R}{\mathbb{R}}
\newcommand{\brac}[1]{\left( #1\right)}
\newcommand{\norm}[1]{\left\lVert #1 \right\rVert}
\newcommand{\proj}{\operatorname{proj}}
\newcommand{\vol}{\mathsf{vol}}
\title{Oracle-based Uniform Sampling from Convex Bodies}
\author{Thanh Dang, Jiaming Liang}
\date{Source: arXiv:2510.02983v2 (CC BY 4.0)}
\begin{document}
\maketitle

\noindent\emph{Source and licence.} Oracle-based Uniform Sampling from Convex Bodies. Version arXiv:2510.02983v2 (CC BY 4.0), distributed by arXiv under the Creative Commons Attribution 4.0 International licence, http://creativecommons.org/licenses/by/4.0/, as recorded in the arXiv licence metadata for this exact version and shown on the abstract page https://arxiv.org/abs/2510.02983v2. Full licence notice: \texttt{licenses/arxiv-2510.02983-CC-BY-4.0.txt}.\par\smallskip

\noindent\textit{Editorial scope.} Counted unit: the article's lemma lem:alternativeintbypart (source0.tex lines 956--974). The article's complete original proof of it is delivered with it (source0.tex lines 977--1066, one \texttt{proof} environment of 90 lines). No mathematics is written, repaired or re-derived: the statement and the proof are the article's own lines, in the article's own notation, and no retained line is substituted or re-flowed.  Retained uncounted as the source-local context the proof needs: lines 418--424; lines 748--757.  Nothing was omitted from any retained span, so the package carries no omission record.  No retained line contains a citation, so the wrapper carries no bibliography and no bibliography entry.  No retained line contains a figure environment, an image-inclusion command or drawing code, so no image is delivered and the package declares no asset dependency.  Editorial formatting only, in addition: an AMS \texttt{amsart} preamble that restates the article's own declarations and macros used by the retained lines, namely the environments \texttt{lemma} on separate counters, together with the standard packages those lines need.  The retained environments print in this document as Lemma 1 in order of appearance, whereas the article prints the same items with the numbering of its own full document.  The complete native source of the article as distributed in the arXiv e-print for this exact version is bundled unchanged as \texttt{sources/186326/source0.tex}.
\begin{lemma}\label{lem:intR}
Let $\tau\ge 0$ be given and assume condition (A1) holds, then we have
\[
      \int_{\R^d} \exp\brac{-\frac{\brac{\|\proj_K(y)-y\|-\tau}^2}{2\eta}}\D y 
    \le \vol(K) \left[\exp\left(\frac{\eta d^2}{2}+\tau d\right) \sqrt{2\pi \eta d^2}+\exp\brac{-\frac{\tau^2}{2\eta}}\right]. 
\]
\end{lemma}

\begin{lemma}\label{lem:gaussianint}
The following statements hold for any $\eta>0$, $c\in \R^d$ and $b\in \R$.
\begin{itemize}
    \item[(a)] $ \int_{\R^d} \exp\brac{-\frac{1}{2\eta} \norm{x-c}^2} \D x=(2\pi \eta)^{d/2}$;

    \item[(b)] $ \int_{\R^d} \exp\left(-\frac{1}{2\eta}(\|x-c\|-b)^2\right) \D x
    \le \exp \left(\frac14 + \frac{d}{\eta}b^2\right)(2\pi\eta)^{d/2}.$

\end{itemize} 
\end{lemma}

\begin{lemma}
\label{lem:alternativeintbypart}
Assume the setup in Lemma~\ref{lem:intR} and write $v(\delta)= \mathrm{vol}(K_\delta)$ and $a(\delta) = \mathrm{vol}_{d-1}(\partial K_\delta)$. Then, $v'(\delta)=a(\delta)$ for almost every
$\delta\ge 0$.
Fix $\eta>0$ and $\tau\ge 0$, and set
\[
  w_{\eta,\tau}(\delta)
  = \exp\!\left(-\frac{(\delta-\tau)^2}{2\eta}\right).
\]
It holds that
\[
\int_0^\infty a(\delta) w_{\eta,\tau}(\delta) \D\delta
\le
\vol(K)\Bigl(e^{d\tau}-1\Bigr)
+
\vol(K)\sqrt{2\pi\eta d^2}\;
\exp\!\left(\tau d+\frac{\eta d^2}{2}\right).
\]
\end{lemma}

\begin{proof}
By the co-area formula, we have $v(\delta) = \mathrm{vol}(K_\delta) = \int_0^{\delta} \mathrm{vol}_{d-1}(\partial K_s) \D s = \int_0^{\delta} a(s) \D s$, so that  $v'(\delta)=a(\delta)$ for almost every $\delta\ge 0$. Then integrating by parts gives
\begin{align}
\label{intbypart}
 \int_0^\infty a(\delta) w_{\eta,\tau}(\delta)  \D \delta=  \int_0^\infty v'(\delta) w_{\eta,\tau}(\delta)  \D \delta
  =
  v(\delta)w_{\eta,\tau}(\delta)\bigg|_{0}^{\infty}
  - \int_0^\infty v(\delta) w_{\eta,\tau}'(\delta)  \D \delta.
\end{align}
Noting that $v(\delta)w_{\eta,\tau}(\delta)\to 0$ as $\delta\to\infty$, we have
\begin{equation}
\label{step_intbypart}
  \int_0^\infty a(\delta) w_{\eta,\tau}(\delta)  \D \delta
  =
  -\mathrm{vol} (K)\exp\!\left(-\frac{\tau^2}{2\eta}\right)
  + \int_0^\infty v(\delta) \frac{\delta-\tau}{\eta} w_{\eta,\tau}(\delta)  \D \delta.
\end{equation}
It follows from the observation $K\subseteq K_\delta\subseteq (1+\delta)K$ that
\begin{equation}
\label{eq:vdeltacomparedtov0}
  \vol(K)\le v(\delta)\le (1+\delta)^d \mathrm{vol} (K)\le e^{ d \delta} \mathrm{vol} (K).
\end{equation}
For $\delta \in [0,\tau]$, we have
\[
\int_0^\tau v(\delta)\frac{\delta-\tau}{\eta}w_{\eta,\tau}(\delta) \D\delta
\stackrel{\eqref{eq:vdeltacomparedtov0}}\le
\vol(K)\int_0^\tau \frac{\delta-\tau}{\eta}w_{\eta,\tau}(\delta) \D\delta
=
\vol(K)\Bigl(-1+e^{-\tau^2/(2\eta)}\Bigr),
\]
where the last identity uses $w'_{\eta,\tau}(\delta)=-\frac{\delta-\tau}{\eta}w_{\eta,\tau}(\delta)$.
For $\delta \in [\tau,\infty)$, using \eqref{eq:vdeltacomparedtov0}, we obtain
\[
\int_\tau^\infty v(\delta)\frac{\delta-\tau}{\eta}w_{\eta,\tau}(\delta) \D\delta
\le
\vol(K)\int_\tau^\infty e^{d\delta}\frac{\delta-\tau}{\eta}w_{\eta,\tau}(\delta) \D\delta.
\]
Plugging the above two inequalities into \eqref{step_intbypart} yields
\begin{equation}
\int_0^\infty a(\delta) w_{\eta,\tau}(\delta) \D\delta
=
-\vol(K)
+\vol(K)\int_\tau^\infty e^{d\delta}\frac{\delta-\tau}{\eta}w_{\eta,\tau}(\delta) \D\delta.
\label{step_intbypart_split}
\end{equation}
Let $F(\delta):=e^{d\delta}w_{\eta,\tau}(\delta)$, then we note that
\[
  \frac{\delta-\tau}{\eta}F(\delta)=dF(\delta)-F'(\delta).
\]
It follows from $F(\delta)\to 0$ as $\delta\to\infty$ and $F(\tau)=e^{d\tau}$ that
\begin{align}
\int_\tau^\infty e^{d\delta}\frac{\delta-\tau}{\eta}w_{\eta,\tau}(\delta) \D\delta
&=
\int_\tau^\infty \Bigl(dF(\delta)-F'(\delta)\Bigr) \D\delta \nonumber\\
&=
d\int_\tau^\infty e^{d\delta}w_{\eta,\tau}(\delta) \D\delta
-F(\delta)\bigg|_{\tau}^{\infty}
=
d\int_\tau^\infty e^{d\delta}w_{\eta,\tau}(\delta) \D\delta
+e^{d\tau}.
\label{eq:kindofibp_alt}
\end{align}
Plugging \eqref{eq:kindofibp_alt} into \eqref{step_intbypart_split} yields
\begin{align}
\int_0^\infty a(\delta) w_{\eta,\tau}(\delta) \D\delta
\le
\vol(K)\Bigl(e^{d\tau}-1\Bigr)
+
d \vol(K)\int_\tau^\infty e^{d\delta}w_{\eta,\tau}(\delta) \D\delta.
\label{step_reduce_to_gaussian_tail}
\end{align}
Finally, we complete the square by writing $d\delta-\frac{(\delta-\tau)^2}{2\eta}
=
-\frac{\left(\delta-(\tau+\eta d)\right)^2}{2\eta}
+\tau d+\frac{\eta d^2}{2}$,
so
\begin{align}
\int_\tau^\infty e^{d\delta}w_{\eta,\tau}(\delta) \D\delta
&\le
\int_{-\infty}^{\infty}\exp\!\left(d\delta-\frac{(\delta-\tau)^2}{2\eta}\right) \D\delta \nonumber\\
&=
\exp\!\left(\tau d+\frac{\eta d^2}{2}\right)
\int_{-\infty}^{\infty}\exp\!\left(-\frac{\left(\delta-(\tau+\eta d)\right)^2}{2\eta}\right) \D\delta \nonumber\\
&=
\sqrt{2\pi\eta}\;
\exp\!\left(\tau d+\frac{\eta d^2}{2}\right),
\end{align}
where the last identity follows from Lemma~\ref{lem:gaussianint}(a).
Combining this with \eqref{step_reduce_to_gaussian_tail} gives, for every $\tau\ge 0$, the desired bound on $\int_0^\infty a(\delta) w_{\eta,\tau}(\delta) \D\delta$. This completes the proof.
\end{proof}

\end{document}
